\documentclass[10pt,a4paper]{amsart}
\usepackage{amsmath,amssymb,amsthm,graphicx}
% For commutative diagrams we use Xy-pic
\usepackage[all]{xy}

% We use 2cell for 2-commutative diagrams.
\xyoption{2cell}
% Theorem environments.
%
\theoremstyle{plain}
\newtheorem{theorem}[subsection]{Theorem}
\newtheorem{proposition}[subsection]{Proposition}
\newtheorem{lemma}[subsection]{Lemma}

\theoremstyle{definition}
\newtheorem{definition}[subsection]{Definition}
\newtheorem{example}[subsection]{Example}
\newtheorem{exercise}[subsection]{Exercise}
\newtheorem{situation}[subsection]{Situation}

\theoremstyle{remark}
\newtheorem{remark}[subsection]{Remark}
\newtheorem{remarks}[subsection]{Remarks}

\numberwithin{equation}{subsection}

% Macros
%
\def\lim{\mathop{\mathrm{lim}}\nolimits}
\def\colim{\mathop{\mathrm{colim}}\nolimits}
\def\Spec{\mathop{\mathrm{Spec}}}
\def\Hom{\mathop{\mathrm{Hom}}\nolimits}
\def\Ext{\mathop{\mathrm{Ext}}\nolimits}
\def\SheafHom{\mathop{\mathcal{H}\!\mathit{om}}\nolimits}
\def\SheafExt{\mathop{\mathcal{E}\!\mathit{xt}}\nolimits}
\def\Sch{\mathit{Sch}}
\def\Mor{\mathop{\mathrm{Mor}}\nolimits}
\def\Ob{\mathop{\mathrm{Ob}}\nolimits}
\def\Sh{\mathop{\mathit{Sh}}\nolimits}
\def\NL{\mathop{N\!L}\nolimits}
\def\CH{\mathop{\mathrm{CH}}\nolimits}
\def\proetale{{pro\text{-}\acute{e}tale}}
\def\etale{{\acute{e}tale}}
\def\QCoh{\mathit{QCoh}}
\def\Ker{\mathop{\mathrm{Ker}}}
\def\Im{\mathop{\mathrm{Im}}}
\def\Coker{\mathop{\mathrm{Coker}}}
\def\Coim{\mathop{\mathrm{Coim}}}

% Boxtimes
%
\DeclareMathSymbol{\boxtimes}{\mathbin}{AMSa}{"02}

%
% Macros for moduli stacks/spaces
%
\def\QCohstack{\mathcal{QC}\!\mathit{oh}}
\def\Cohstack{\mathcal{C}\!\mathit{oh}}
\def\Spacesstack{\mathcal{S}\!\mathit{paces}}
\def\Quotfunctor{\mathrm{Quot}}
\def\Hilbfunctor{\mathrm{Hilb}}
\def\Curvesstack{\mathcal{C}\!\mathit{urves}}
\def\Polarizedstack{\mathcal{P}\!\mathit{olarized}}
\def\Complexesstack{\mathcal{C}\!\mathit{omplexes}}
% \Pic is the operator that assigns to X its picard group, usage \Pic(X)
% \Picardstack_{X/B} denotes the Picard stack of X over B
% \Picardfunctor_{X/B} denotes the Picard functor of X over B
\def\Pic{\mathop{\mathrm{Pic}}\nolimits}
\def\Picardstack{\mathcal{P}\!\mathit{ic}}
\def\Picardfunctor{\mathrm{Pic}}
\def\Deformationcategory{\mathcal{D}\!\mathit{ef}}
\usepackage{cleveref}
\crefname{originalreference}{original reference}{original references}
\Crefname{originalreference}{Original reference}{Original references}
\makeatletter
\newlabel{situation-setup}{{02QL}{source}{}{sourcepointer-1}{}}
\newlabel{situation-setup@cref}{{[originalreference][0][]02QL}{[1][0][]source}{}{}{}}
\makeatother
\crefname{originallemma}{original lemma}{original lemmas}
\Crefname{originallemma}{Original lemma}{Original lemmas}
\makeatletter
\newlabel{algebra-lemma-reformulate-CM}{{02JN}{source}{}{sourcepointer-2}{}}
\newlabel{algebra-lemma-reformulate-CM@cref}{{[originallemma][0][]02JN}{[1][0][]source}{}{}{}}
\makeatother
\crefname{originallemma}{original lemma}{original lemmas}
\Crefname{originallemma}{Original lemma}{Original lemmas}
\makeatletter
\newlabel{divisors-lemma-Noetherian-scheme-regular-ideal}{{063I}{source}{}{sourcepointer-3}{}}
\newlabel{divisors-lemma-Noetherian-scheme-regular-ideal@cref}{{[originallemma][0][]063I}{[1][0][]source}{}{}{}}
\makeatother
\crefname{originallemma}{original lemma}{original lemmas}
\Crefname{originallemma}{Original lemma}{Original lemmas}
\makeatletter
\newlabel{lemma-gysin-fundamental}{{0FBM}{source}{}{sourcepointer-4}{}}
\newlabel{lemma-gysin-fundamental@cref}{{[originallemma][0][]0FBM}{[1][0][]source}{}{}{}}
\makeatother
\title{The Gysin class of a virtual normal sheaf of rank r caps the fundamental class of a Cohen-Macaulay situation to the fundamental class of the intersection}
\author{The Stacks Project Authors}
\date{stacks-project:chow.tex:local-snapshot-20261012; GNU FDL 1.2 or later}
\begin{document}
\maketitle
\noindent\textit{Editorial note.} All delivered mathematical text is the author's own, copied from the pinned Stacks Project chapter file chow.tex; no statement or proof step was rewritten, reordered or removed.\par\medskip
\noindent\textit{Editorial note.} Preamble provenance: the pinned chapter file opens with an input of the shared preamble and carries no preamble text of its own, so the project's single Stacks-wide preamble staged at sources/stacks-chow/preamble.tex is excerpted verbatim as two spans, lines 14-18 (Xy-pic setup) and lines 154-217 (the native theorem declarations and mathematical macros). Nothing was added, removed or redefined; the wrapper class and the amsmath/amssymb/amsthm/graphicx base packages are supplied by the assembler.\par\medskip
\noindent\textit{Editorial note.} The counted statement and proof contain no cross reference, no citation and no included file, so the delivered body is self contained.\par\medskip
\noindent\textit{Editorial note.} Canonical location of the counted proposition: Stacks tag 0FEA (sources/stacks-chow/pinned-original-source.tex line 12128).\par\medskip
\noindent\textit{Allowed standard prerequisites (editorial).}\begin{itemize}
\item Cohen-Macaulay local rings are exactly those whose ideal of definition is generated by a regular sequence. Algebra, Lemma labelled algebra-lemma-reformulate-CM (Stacks tag 02JN), invoked at proof lines 12158-12162. Supplies that I\_xi is generated by a regular sequence, which is the local regularity input.
\item A locally generated-by-r-elements ideal whose localisation is regular-sequence generated is so on a neighbourhood. Divisors, Lemma labelled divisors-lemma-Noetherian-scheme-regular-ideal (Stacks tag 063I), invoked at proof lines 12163-12164. Spreads the regular sequence from the generic point to an open neighbourhood V of it.
\item The fundamental-class identity for a Gysin class with vanishing excess normal sheaf. Same chapter, Lemma labelled lemma-gysin-fundamental (source line 12076), invoked at proof lines 12168-12169. Supplies the conclusion once the excess normal sheaf is zero over V.
\item The base situation (S, delta) for operational Chow groups and virtual normal sheaves. Same chapter, Situation labelled situation-setup (source line 1372), invoked in the statement. Fixes the dimension function and the ring A\textasciicircum{}*(X) in which the Gysin class lives.
\end{itemize}\medskip
\ifdefined\hypertarget\hypertarget{sourcepointer-1}{}\fi
\noindent\textit{Source reference.} Original reference 02QL: sources/stacks-chow/pinned-original-source.tex, source lines 1372--1372. Same-chapter pointer into the author environment situation at source line 1372; invoked at proof line None of the counted proof. It is retained only as a navigation pointer; every load-bearing step is inside the delivered span.\par\medskip
\ifdefined\hypertarget\hypertarget{sourcepointer-2}{}\fi
\noindent\textit{Source reference.} Original lemma 02JN: sources/stacks-chow/tags.txt, source lines 3342--3342. Cross-chapter standard prerequisite located in the project tag list staged in the local snapshot; invoked at proof line 12162 of the counted proof. It is retained only as a navigation pointer; every load-bearing step is inside the delivered span.\par\medskip
\ifdefined\hypertarget\hypertarget{sourcepointer-3}{}\fi
\noindent\textit{Source reference.} Original lemma 063I: sources/stacks-chow/tags.txt, source lines 7832--7832. Cross-chapter standard prerequisite located in the project tag list staged in the local snapshot; invoked at proof line 12163 of the counted proof. It is retained only as a navigation pointer; every load-bearing step is inside the delivered span.\par\medskip
\ifdefined\hypertarget\hypertarget{sourcepointer-4}{}\fi
\noindent\textit{Source reference.} Original lemma 0FBM: sources/stacks-chow/pinned-original-source.tex, source lines 12076--12076. Same-chapter pointer into the author environment lemma at source line 12076; invoked at proof line 12169 of the counted proof. It is retained only as a navigation pointer; every load-bearing step is inside the delivered span.\par\medskip
\begin{lemma}
\label{lemma-gysin-easy}
Let $(S, \delta)$ be as in Situation \ref{situation-setup}. Let $X$ be a scheme
locally of finite type over $S$. Let $\mathcal{N}$ be a virtual normal
sheaf for a closed subscheme $Z$ of $X$. Let $Y \to X$ be a morphism
which is locally of finite type. Given integers $r$, $n$ assume
\begin{enumerate}
\item $\mathcal{N}$ is locally free of rank $r$,
\item every irreducible component of $Y$ has $\delta$-dimension $n$,
\item $\dim_\delta(Z \times_X Y) \leq n - r$, and
\item for $\xi \in Z \times_X Y$ with $\delta(\xi) = n - r$
the local ring $\mathcal{O}_{Y, \xi}$ is Cohen-Macaulay.
\end{enumerate}
Then $c(Z \to X, \mathcal{N}) \cap [Y]_n = [Z \times_X Y]_{n - r}$
in $\CH_{n - r}(Z \times_X Y)$.
\end{lemma}

\begin{proof}
The statement makes sense as $Z \times_X Y$ is a closed subscheme of $Y$.
Because $\mathcal{N}$ has rank $r$ we know that
$c(Z \to X, \mathcal{N}) \cap [Y]_n$ is in $\CH_{n - r}(Z \times_X Y)$.
Since $\dim_\delta(Z \cap Y) \leq n - r$ the chow group
$\CH_{n - r}(Z \times_X Y)$ is freely generated by the
cycle classes of the irreducible components $W \subset Z \times_X Y$
of $\delta$-dimension $n - r$. Let $\xi \in W$ be the generic point.
By assumption (2) we see that $\dim(\mathcal{O}_{Y, \xi}) = r$.
On the other hand, since $\mathcal{N}$ has rank $r$ and since
$\mathcal{N}^\vee \to \mathcal{C}_{Z/X}$ is surjective, we see that
the ideal sheaf of $Z$ is locally cut out by $r$ equations.
Hence the quasi-coherent ideal sheaf $\mathcal{I} \subset \mathcal{O}_Y$
of $Z \times_X Y$ in $Y$ is locally generated by $r$ elements.
Since $\mathcal{O}_{Y, \xi}$ is Cohen-Macaulay of dimension $r$
and since $\mathcal{I}_\xi$ is an ideal of definition (as $\xi$ is
a generic point of $Z \times_X Y$) it follows that $\mathcal{I}_\xi$
is generated by a regular sequence
(Algebra, Lemma \ref{algebra-lemma-reformulate-CM}).
By Divisors, Lemma \ref{divisors-lemma-Noetherian-scheme-regular-ideal}
we see that $\mathcal{I}$ is generated by a regular sequence over
an open neighbourhood $V \subset Y$ of $\xi$. By our description of
$\CH_{n - r}(Z \times_X Y)$ it suffices to show that
$c(Z \to X, \mathcal{N}) \cap [V]_n = [Z \times_X V]_{n - r}$
in $\CH_{n - r}(Z \times_X V)$. This follows from
Lemma \ref{lemma-gysin-fundamental}
because the excess normal sheaf is $0$ over $V$.
\end{proof}

\end{document}
